\chapter{Нулевой порог канонического дефекта и канал полного момента один}

Предыдущая глава не нашла отрицательных мод, но нижняя тройка собственных
значений уменьшалась при расширении радиального базиса. Теперь необходимо
решить, является ли это началом неустойчивости или приближением к
неотрицательному непрерывному спектру безмассовых мод.

\section{Литературная и проектная граница}

Для солитонов нулевые моды связаны с нарушенными непрерывными симметриями и
должны отделяться от внутренних колебаний; систематический принцип
коллективных координат изложен в \path{arXiv:1008.3605}. В моделях Скирма
три переноса и три вращения являются стандартными нулевыми направлениями;
см. \path{arXiv:hep-ph/9211229}. Частично-волновой анализ явно помещает
переносную моду ежа в канал момента один; пример дан в
\path{arXiv:hep-ph/9907492}.

Эти результаты не определяют спектр нашего составного $Q$-дефекта. Они
задают правильный вопрос: сначала установить представление нижней моды и
её масштабирование, а уже затем решать, является ли она нормируемой
коллективной координатой, пороговым резонансом или отрицательной модой.

\section{Диагональный полный угловой момент}

Канонический ёж инвариантен не относительно независимых пространственных и
семейных вращений, а относительно их диагональной комбинации. Поле $Q$
несёт внутренний спин два. Поэтому скалярная орбитальная гармоника степени
$\ell$ раскладывается по полному моменту $J$ как
\begin{equation}
 \ell\otimes2
 =\bigoplus_{J=|\ell-2|}^{\ell+2}J.
 \label{eq:v6-grand-angular-momentum-decomposition}
\end{equation}
В частности,
\begin{equation}
 2\otimes2=0\oplus1\oplus2\oplus3\oplus4,
\end{equation}
с размерностями $1,3,5,7,9$.

Нижняя мода предыдущего аудита целиком лежала в скалярном секторе
$\ell=2$ и была трёхкратно вырождена. В разложении
\eqref{eq:v6-grand-angular-momentum-decomposition} единственное
трёхмерное неприводимое представление имеет $J=1$.

\begin{proposition}[Классификация нижней тройки]
Нижняя трёхкратно вырожденная мода пятикомпонентного гессиана принадлежит
каналу диагонального полного углового момента $J=1$.
\end{proposition}

Это тот же симметрийный тип, что и у переносов ежа, но одного совпадения
представлений недостаточно для объявления моды нормируемой коллективной
координатой.

\section{Масштабирование к бесконечному объёму}

Для выделенного $\ell=2$-сектора построены пространства размерности $125$
на шарах радиусов
\begin{equation}
 L=6,8,10,12.
\end{equation}
Пять радиальных масштабов включали две функции ядра и три функции, размеры
которых росли пропорционально $L$. Нижние собственные значения равны
\begin{equation}
 \lambda_1(L)=
 \{0.05050663\ldots,
 0.01374853\ldots,
 0.00536173\ldots,
 0.00275369\ldots\}.
 \label{eq:v6-J1-finite-volume-sequence}
\end{equation}
Во всех четырёх случаях вся нижняя тройка остаётся положительной.

Степенная аппроксимация последних трёх точек даёт
\begin{equation}
 \lambda_1(L)\sim L^{-3.97514\ldots}.
 \label{eq:v6-J1-quartic-threshold-fit}
\end{equation}
Произведения $L^4\lambda_1$ равны
$56.31\ldots,53.62\ldots,57.10\ldots$ для $L=8,10,12$.

\section{Почему возникает четвёртая степень}

На упорядоченной вакуумной орбите
\begin{equation}
 Q=\Delta\left(P-\frac{I_3}{3}\right),
 \qquad A_Q=[P,dP].
\end{equation}
Проекторное тождество даёт
\begin{equation}
 D_QQ=dQ+[A_Q,Q]=0,
 \qquad V_{\rm can}(Q)=0.
 \label{eq:v6-vacuum-orbit-two-derivative-cancellation}
\end{equation}
Следовательно, длинноволновое изменение директора не имеет ни объёмной,
ни двухпроизводной цены. Первый ненулевой статический вклад имеет вид
\begin{equation}
 \int|F_{A_Q}|^2,
\end{equation}
и содержит четыре производные в квадратичной длинноволновой теории.
При нормировке $\int|\delta Q|^2=1$ это предсказывает
$\lambda\sim k^4\sim L^{-4}$, в согласии с
\eqref{eq:v6-J1-quartic-threshold-fit}.

\begin{theorem}[Статический нулевой порог]
Нижняя $J=1$-тройка в проверенной последовательности конечных объёмов
приближается к нулю сверху с законом $L^{-4}$. Это поведение следует из
точного обнуления двухпроизводного и потенциального вкладов на вакуумной
орбите и не является численным признаком отрицательной моды.
\end{theorem}

\section{Первый старший угловой сектор}

Для $\ell=3$ добавлены семь регулярных гармонических полиномов и четыре
радиальные функции. Размер пятикомпонентного пространства равен $140$.
Нижнее собственное значение получилось
\begin{equation}
 \lambda_{\ell=3}^{\min}=2.82040280\ldots,
\end{equation}
отрицательных мод не найдено, а проекция первой вариации имеет норму
$5.9\cdot10^{-12}$.

Этот тест исключает немедленное появление неустойчивости в первом ранее не
проверенном угловом секторе. Он не является доказательством для всех
$\ell\ge3$.

\section{Нормировка и физическая граница}

Знак статической квадратичной формы не зависит от выбора положительной
нормы. Однако численные значения $\lambda$ и закон
\eqref{eq:v6-J1-quartic-threshold-fit} записаны относительно
$L^2$-нормы поля $Q$. Текущий родитель всё ещё не фиксирует единственную
лоренцеву кинетическую метрику, включая временную компоненту составной
связности. Поэтому из статического гессиана нельзя извлекать физические
частоты или скорость распространения мод.

\begin{nogocriterion}[Порог не равен полной теореме устойчивости]
Идентификация $J=1$-порога и положительность проверенного $\ell=3$-сектора
не доказывают неотрицательность всех радиальных операторов при каждом $J$.
Следующая проверка должна вывести канальные операторы и попытаться
представить их как сумму положительных квадратов и конечного матричного
потенциала.
\end{nogocriterion}

\begin{protocol}
\begin{enumerate}
 \item внутренний спин поля $Q$: \textbf{$2$};
 \item симметрия ежа: \textbf{диагональный полный момент};
 \item нижняя тройка: \textbf{$J=1$};
 \item радиусы конечных областей: \textbf{$6,8,10,12$};
 \item знак нижней тройки: \textbf{положительный во всех случаях};
 \item асимптотическая степень: \textbf{$-3.97514\ldots$};
 \item аналитическое ожидание: \textbf{$L^{-4}$};
 \item статический непрерывный спектр начинается с нуля:
 \textbf{подтверждено};
 \item отрицательное связанное состояние: \textbf{не найдено};
 \item сектор $\ell=3$: \textbf{положителен в проверенном базисе};
 \item физические частоты: \textbf{не выведены};
 \item полная операторная неотрицательность: \textbf{ещё не доказана};
 \item следующий гейт: \textbf{факторизация канальных операторов}.
\end{enumerate}
\end{protocol}

Машинный сертификат сохранён в
\path{s2t_v6_bosonic_defect_canonical_continuum_stability_gate_results.json}.